\documentclass[11pt]{article}

\usepackage[margin=1in]{geometry}
\usepackage{graphicx}
\usepackage{booktabs}
\usepackage{caption}
\usepackage{subcaption}
\usepackage{amsmath}
\usepackage{amssymb}
\usepackage{hyperref}
\usepackage{xcolor}
\usepackage{placeins}
\usepackage{array}
\setlength{\emergencystretch}{2em}
\newcolumntype{L}[1]{>{\raggedright\arraybackslash}p{#1}}
\newcommand{\ind}{\mathbf{1}}
\DeclareMathOperator{\logit}{logit}

\hypersetup{
  colorlinks=true,
  linkcolor=black,
  citecolor=black,
  urlcolor=blue
}

\graphicspath{{../outputs/figures/}}

\title{A Public-Data Bayesian Framework for Retrospective External Wildfire Exposure Ranking of Transmission-Line Segments}
\author{Wenxing Yi}
\date{July 7, 2026 revision}

\begin{document}
\maketitle

\begin{abstract}
Transmission-line wildfire decisions often require segment-level triage, while
equipment-condition, outage, and inspection records are not public. This paper
presents a reproducible public-data framework for retrospective external
wildfire exposure ranking of California transmission-line segments. The line
segments are treated only as receptor assets. They are not modeled as ignition
sources, failure points, outage causes, or responsibility units. Using
California Energy Commission transmission lines, CAL FIRE FRAP fire perimeters
and cause fields, gridMET weather summaries, LANDFIRE LF2024 fuel and canopy
layers, and USGS 3DEP terrain, the workflow constructs a 2017--2023
segment-year panel for the Northern Sierra and Southern Cascades. The main
label excludes CAL FIRE CAUSE = 11, Electrical Power, to reduce conflation
between external wildfire exposure and electrical-power fire cause. A Bayesian
logistic exposure model is compared with a deterministic public physics score
and machine-learning baselines under a 2022--2023 temporal holdout. Bayesian
Laplace and L2 logistic discrimination are close; the Bayesian contribution is
therefore framed as posterior exceedance probability and posterior top-decile
rank probability for rare-event triage, not as universal classification
superiority. Probability calibration remains limited relative to the rare-event
climatology baseline, so posterior rank probability is treated as the primary
decision output.
\end{abstract}

\section{Introduction}

Wildfire exposure increasingly affects linear infrastructure systems. For
transmission corridors, screening decisions are often made at the segment scale
with incomplete public information about asset condition, protection systems,
inspection history, vegetation work, and outage outcomes. A public-data
workflow can therefore support transparent exposure ranking, but it should not
be presented as a full operational risk model.

This study is intentionally narrow. It addresses fire-to-grid external exposure
ranking for transmission-line receptor assets. It does not estimate
powerline-caused ignition probability, equipment failure, outage probability,
damage, or legal responsibility. Electrical-power fire overlap is retained only
as a diagnostic, because a perimeter-level cause field does not identify the
exact ignition device or segment.

The methodological question is also narrow. Deterministic public physics scores
and machine-learning classifiers return useful point scores, but rare-event
triage often asks threshold and ranking questions: which segments plausibly
belong in the high-exposure tail, and how uncertain is that ranking? A Bayesian
decision layer answers these questions with posterior exceedance probability
and posterior top-decile rank probability. This is the main contribution, rather
than a claim that Bayesian point prediction dominates all alternatives.

The contributions are:
\begin{enumerate}
  \item a reproducible public-data segment-year exposure panel for California
  transmission-line receptor assets;
  \item cause-screened external exposure labels that exclude electrical-power
  fire causes from the main validation target;
  \item weather, terrain, and fuel-structure covariates for retrospective
  annual exposure ranking;
  \item Bayesian posterior exceedance and posterior top-decile rank
  probabilities as decision-facing outputs; and
  \item temporal-holdout comparison with deterministic, L2 logistic, tree-based
  machine-learning, and initial Stan MCMC baselines.
\end{enumerate}

\section{Study Area and Public Data}

\subsection{Study Area}

The case study focuses on the Northern Sierra and Southern Cascades in
California. Decision units are approximately 1 km transmission-line segments.
The analysis panel spans 2017--2023 and uses segment-year rows so annual fire
labels and annual observed environmental summaries are aligned. The current
prototype uses a fixed 1 km segment buffer for perimeter-overlap labels; buffer
sensitivity is a required follow-up analysis. \autoref{fig:study-area} shows
the study region and receptor-asset segments.

\begin{figure}[!htbp]
  \centering
  \includegraphics[width=0.94\textwidth]{sci_phase1/fig1_study_area_transmission_fire_labels.png}
  \caption{Study region, public transmission-line segments, and annual exposure labels.}
  \label{fig:study-area}
\end{figure}

\subsection{Public Data Sources}

\autoref{tab:data-sources} summarizes the public data sources. The processed
panel contains 11,255 line segments and 78,785 segment-year observations. The
environmental covariate maps are treated as covariate diagnostics and are moved
to the Appendix so that the main paper emphasizes label construction,
validation, and decision quantities.

\begin{table}[!htbp]
  \centering
  \footnotesize
  \setlength{\tabcolsep}{3pt}
  \caption{Public data sources and analytical roles.}
  \label{tab:data-sources}
  \begin{tabular}{L{0.16\textwidth} L{0.18\textwidth} L{0.16\textwidth} L{0.14\textwidth} L{0.25\textwidth}}
    \toprule
    Dataset & Source & Resolution / unit & Version or years & Role and caveat \\
    \midrule
    Transmission lines & California Energy Commission & Public line geometry & Current public release & Receptor assets only; not equipment-failure data. \\
    Fire perimeters & CAL FIRE FRAP & Perimeter polygons with cause fields & 2017--2023 subset & Exposure labels; not damage or outage observations. \\
    gridMET & University of Idaho gridMET & Daily weather, annual summaries & 2017--2023 & Observed annual weather covariates. \\
    LANDFIRE & LANDFIRE LF2024 & Fuel and canopy rasters & LF2024 & Static fuel-structure covariates; time-matched sensitivity is required. \\
    USGS 3DEP & USGS 3DEP & Elevation raster & Current public DEM & Terrain covariate; not a fire-spread model. \\
    \bottomrule
  \end{tabular}
\end{table}

\section{Label Construction}

Let $i$ index line segments, $t$ years, and $j$ individual fire perimeters.
Let $B_i$ be the segment buffer and $P_j$ be fire perimeter $j$. The main
external exposure label is
\begin{equation}
  y^{exo}_{it}
  =
  \ind\left\{
  \exists j:\; \mathrm{year}(j)=t,\;
  B_i \cap P_j \neq \emptyset,\;
  \mathrm{CAUSE}_j \neq 11
  \right\}.
  \label{eq:y-exo}
\end{equation}
CAL FIRE CAUSE = 11 denotes Electrical Power. The exclusion in
\autoref{eq:y-exo} is used to avoid mixing external exposure with
electrical-power fire cause in the main validation target. It does not assert
that remaining fires cannot affect or be affected by transmission assets.

Four annual labels are retained (\autoref{tab:label-algorithm}). The all-fire
label records any perimeter overlap. The exogenous label is the main target.
The strict label further excludes Equipment Use. Electrical-power overlap is a
diagnostic only and is not interpreted as ignition probability.

\begin{table}[!htbp]
  \centering
  \scriptsize
  \setlength{\tabcolsep}{2pt}
  \caption{Annual segment-year labels. Labels are not mutually exclusive because multiple perimeters can overlap a segment-year.}
  \label{tab:label-algorithm}
  \begin{tabular}{L{0.20\textwidth} L{0.42\textwidth} r L{0.22\textwidth}}
    \toprule
    Label & Definition & Positives & Interpretation \\
    \midrule
    All-fire exposure & Any 1 km segment-buffer intersection with a CAL FIRE perimeter & 3,187 & Broad exposure screen. \\
    Exogenous exposure & Any overlap after excluding CAUSE = 11 & 2,331 & Main external exposure target. \\
    Strict exogenous exposure & Excludes CAUSE in \{Electrical Power, Equipment Use\} & 2,075 & Sensitivity label. \\
    Electrical-power overlap & Any overlap with CAUSE = 11 & 885 & Diagnostic only. \\
    Mixed exogenous/electrical rows & Segment-years with both exogenous and electrical-power overlaps & 29 & Multiple-fire overlap diagnostic. \\
    \bottomrule
  \end{tabular}
\end{table}

\begin{figure}[!htbp]
  \centering
  \includegraphics[width=0.94\textwidth]{sci_phase1/fig2_cause_screened_label_maps.png}
  \caption{Cause-screened exposure labels used for receptor-asset validation.}
  \label{fig:cause-screening}
\end{figure}

\begin{figure}[!htbp]
  \centering
  \includegraphics[width=0.88\textwidth]{sci_phase1/fig3_annual_segment_year_label_counts.png}
  \caption{Annual segment-year label counts.}
  \label{fig:annual-labels}
\end{figure}
\FloatBarrier

\section{Methods}

\subsection{Feature Core}

The selected feature core combines observed annual gridMET dryness and wind
summaries, USGS terrain, LANDFIRE canopy/fuel-structure variables, and simple
physics-guided indices. In this paper, physics-guided means that covariates are
engineered from known weather, fuel, and terrain controls; it does not mean that
the model simulates fire spread.

Let $z(\cdot)$ denote training-period standardization. The dryness index is
\begin{equation}
  \begin{split}
    d_{it}
    &=
    \frac{z(\mathrm{VPD}_{it})+z(\mathrm{ERC}_{it})+z(\mathrm{BI}_{it})+
    z(\mathrm{wind}_{it})}{4}\\
    &\quad -
    \frac{z(\mathrm{FM100}_{it})+z(\mathrm{FM1000}_{it})+
    z(\mathrm{RMIN}_{it})+z(\mathrm{PR}_{it})}{4}.
  \end{split}
\end{equation}
The fuel-structure index is
\begin{equation}
  f_i =
  \frac{
  z(\mathrm{canopy\ cover}_i)+
  z(\mathrm{canopy\ height}_i)+
  z(\mathrm{canopy\ bulk\ density}_i)
  }{3}.
\end{equation}
The selected core adds $d_{it}f_i$ to the raw public covariates. LANDFIRE
FBFM40 is categorical and is not used in the selected Bayesian or L2 logistic
core. The current Phase 5 ablation scripts used FBFM40 as a continuous
diagnostic variable in several ablations; those ablations are therefore
reported as exploratory until categorical encoding is implemented.

\begin{table}[!htbp]
  \centering
  \small
  \caption{Feature-construction audit for the current revision.}
  \label{tab:feature-audit}
  \begin{tabular}{L{0.26\textwidth} L{0.30\textwidth} L{0.32\textwidth}}
    \toprule
    Component & Current implementation & Revision status \\
    \midrule
    gridMET annual weather & Annual summaries for VPD, ERC, BI, wind, fuel moisture, humidity, and precipitation & Interpreted retrospectively under observed annual conditions. \\
    USGS 3DEP terrain & Segment midpoint elevation and derived terrain covariates & Suitable for screening; not a spread simulation. \\
    LANDFIRE canopy layers & CH/CBH/CBD sampled as raw raster-coded values in processed tables & Report decoded units: CH/CBH raw/10, CBD raw/100. \\
    LANDFIRE FBFM40 & Used as continuous in some exploratory ablations, not in selected core & Must be categorical or removed in journal-grade ablations. \\
    Engineered physics indices & Dryness, fuel structure, dryness-by-fuel interaction & Constant unit rescaling does not change standardized linear fit but does change interpretation. \\
    \bottomrule
  \end{tabular}
\end{table}

\subsection{Deterministic Comparator}

The deterministic comparator is a transparent, calibrated physics-guided score.
It is a benchmark, not the final decision model:
\begin{equation}
  r^{det}_{it}
  =
  0.56 d_{it} + 0.30 f_i + 0.14 z(\mathrm{elevation}_i).
\end{equation}
The raw score is clipped to the training-period 1st--99th percentile range and
scaled to $[0,1]$. It is then calibrated with a regularized logistic model,
\begin{equation}
  p^{det}_{it} = \logit^{-1}(\gamma_0+\gamma_1 s^{det}_{it}).
\end{equation}
This gives a single score and calibrated point probability but no posterior
distribution over ranks.

\subsection{Bayesian and L2 Logistic Models}

The Bayesian exposure model is a weakly regularized logistic regression:
\begin{align}
  y^{exo}_{it} &\sim \mathrm{Bernoulli}(p_{it}), \\
  p_{it} &= \logit^{-1}(\alpha+\mathbf{x}_{it}^{\top}\boldsymbol{\beta}),\\
  \alpha &\sim \mathcal{N}(0,5^2), \qquad
  \beta_k \sim \mathcal{N}(0,2^2).
\end{align}
The posterior is approximated by a Laplace approximation around the posterior
mode $\hat{\theta}$, where
\begin{equation}
  p(\theta \mid D) \approx
  \mathcal{N}\left(\hat{\theta},\; H(\hat{\theta})^{-1}\right),
\end{equation}
and $H(\hat{\theta})$ is the negative Hessian of the log posterior. This is
closely related to L2 logistic regression: a Gaussian prior on coefficients is
equivalent to an L2 penalty at the posterior mode. Thus similar discrimination
between Bayesian Laplace and L2 logistic is expected. The Bayesian layer is
retained because posterior draws support rank and threshold decision quantities.

For posterior draw $s$, segment exposure over holdout years $T^\star$ is
\begin{equation}
  \bar{p}_{i}^{(s)}
  =
  \frac{1}{|T^\star|}
  \sum_{t\in T^\star} \logit^{-1}
  \left(\alpha^{(s)}+\mathbf{x}_{it}^{\top}\boldsymbol{\beta}^{(s)}\right).
\end{equation}
The decision quantities are
\begin{align}
  q_i
  &= P(\bar{p}_i > \tau \mid D)
  \approx
  \frac{1}{S}\sum_{s=1}^{S}
  \ind\{\bar{p}_{i}^{(s)}>\tau\},\\
  \rho_i
  &= P(\mathrm{rank}(\bar{p}_i) \le K \mid D)
  \approx
  \frac{1}{S}\sum_{s=1}^{S}
  \ind\{\mathrm{rank}(\bar{p}_{i}^{(s)})\le K\}.
\end{align}
Here $\tau$ is a prevalence-derived exposure threshold and $K$ is the
top-decile segment count. Because Brier skill is not positive relative to the
holdout climatology baseline, $\rho_i$ is emphasized as the primary decision
quantity and $q_i$ as a secondary threshold diagnostic.

\begin{figure}[!htbp]
  \centering
  \includegraphics[width=0.90\textwidth]{phase8/phase8_bayesian_method_schematic-2.png}
  \caption{Bayesian workflow for posterior exposure and rank probabilities.}
  \label{fig:method-schematic}
\end{figure}

\subsection{Validation Design}

Models are trained on 2017--2021 and evaluated on 2022--2023. The holdout set
contains 22,510 segment-year rows and 356 exogenous exposure positives, giving
a prevalence of 1.5815\%. PR-AUC is therefore interpreted relative to this
rare-event prevalence. Metrics include ROC-AUC, PR-AUC, Brier score, Brier
skill versus climatology, row-level top-k capture, segment-level top-decile
capture, and a preliminary segment-cluster bootstrap. Temporal holdout is not a
substitute for spatial block validation.

\section{Results}

\subsection{Temporal Holdout Performance}

\autoref{tab:main-comparison} summarizes temporal-holdout performance. Bayesian
Laplace, Stan MCMC pilot, and L2 logistic have similar discrimination. Relative
to the deterministic comparator, Bayesian/L2 ranking captures more holdout
positives in the top-ranked tail. Relative to L2 logistic, the Bayesian point
prediction is not materially better; its value is the posterior decision layer.

\begin{table}[!htbp]
  \centering
  \small
  \caption{Temporal-holdout comparison for 2022--2023 exogenous exposure.}
  \label{tab:main-comparison}
  \begin{tabular}{lrrrrr}
    \toprule
    Model & ROC-AUC & PR-AUC & Brier & Row top 10\% & Segment top 10\% \\
    \midrule
    ML logistic L2 & 0.663 & 0.039 & 0.015713 & 0.289 & 0.216 \\
    Bayesian Laplace & 0.662 & 0.038 & 0.015718 & 0.289 & 0.216 \\
    Stan MCMC pilot & 0.662 & 0.038 & 0.015710 & 0.289 & 0.216 \\
    Deterministic calibrated & 0.651 & 0.028 & 0.015734 & 0.219 & 0.118 \\
    Balanced extra trees & 0.617 & 0.022 & 0.056091 & 0.180 & 0.104 \\
    Histogram gradient boosting & 0.553 & 0.019 & 0.016647 & 0.132 & 0.149 \\
    Balanced random forest & 0.568 & 0.018 & 0.034539 & 0.107 & 0.152 \\
    \bottomrule
  \end{tabular}
\end{table}

\begin{figure}[!htbp]
  \centering
  \includegraphics[width=0.88\textwidth]{phase9/phase9_20260707_113214_topk_capture_curve.png}
  \caption{Top-k capture curve under 2022--2023 temporal holdout. Bayesian and L2 logistic rankings are close and both exceed the deterministic comparator in this split.}
  \label{fig:topk-capture}
\end{figure}

\begin{figure}[!htbp]
  \centering
  \includegraphics[width=0.82\textwidth]{phase9/phase9_20260707_113214_brier_skill_baseline.png}
  \caption{Brier skill relative to the holdout climatology baseline. Negative values motivate cautious interpretation of calibrated probabilities.}
  \label{fig:brier-skill}
\end{figure}

\begin{table}[!htbp]
  \centering
  \small
  \caption{Rare-event baseline audit. Holdout prevalence is 1.5815\%, and the climatology Brier score is 0.015565.}
  \label{tab:brier-skill}
  \begin{tabular}{lrrrr}
    \toprule
    Model & PR-AUC & PR-AUC / prevalence & Brier & Brier skill \\
    \midrule
    Logistic L2 & 0.03855 & 2.438 & 0.015713 & -0.0095 \\
    Bayesian Laplace & 0.03838 & 2.427 & 0.015718 & -0.0098 \\
    Deterministic calibrated & 0.02790 & 1.764 & 0.015734 & -0.0109 \\
    \bottomrule
  \end{tabular}
\end{table}

\subsection{Posterior Decision Maps}

The posterior map outputs should be read as uncertainty-aware exposure-ranking
diagnostics. In particular, $\rho_i=P(\mathrm{rank}\le K\mid D)$ identifies
segments that repeatedly fall into the top decile under posterior draws.
\autoref{fig:bayes-maps} and \autoref{fig:posterior-maps} show the posterior
decision metrics for the current prototype.

\begin{figure}[!htbp]
  \centering
  \includegraphics[width=0.94\textwidth]{phase4/phase4_decision_metric_maps.png}
  \caption{Posterior decision metrics mapped to transmission-line segments.}
  \label{fig:bayes-maps}
\end{figure}

\begin{figure}[!htbp]
  \centering
  \includegraphics[width=0.94\textwidth]{phase8/phase8_faceted_zoom_posterior_maps.png}
  \caption{Full-region and zoomed maps of posterior rank and exceedance diagnostics.}
  \label{fig:posterior-maps}
\end{figure}

\subsection{Feature and Method Audits}

\autoref{fig:landfire-audit} shows the LANDFIRE unit audit. The selected
Bayesian and L2 logistic core does not use FBFM40, but several exploratory
ablation models do; those ablations are retained only as diagnostics in this
revision. The current M5 selected core should also be described as a
physics-guided selected core, not as an independently tuned final model unless
pre-specification or nested model selection is documented.

\begin{figure}[!htbp]
  \centering
  \includegraphics[width=0.80\textwidth]{phase9/phase9_20260707_113214_landfire_scaling_audit.png}
  \caption{LANDFIRE unit audit for canopy height, canopy base height, and canopy bulk density.}
  \label{fig:landfire-audit}
\end{figure}

\begin{table}[!htbp]
  \centering
  \small
  \caption{Method audit and current interpretation limits.}
  \label{tab:method-audit}
  \begin{tabular}{L{0.27\textwidth} L{0.31\textwidth} L{0.29\textwidth}}
    \toprule
    Item & Current status & Manuscript interpretation \\
    \midrule
    Temporal holdout & 2017--2021 train, 2022--2023 test & Useful retrospective check; not prospective forecasting proof. \\
    L2 logistic comparison & Nearly identical discrimination to Bayesian Laplace & Bayesian value is posterior decision quantities. \\
    Brier skill & Slightly negative versus climatology & Do not claim calibration superiority. \\
    Segment dependence & Preliminary segment-cluster bootstrap only & Spatial block validation still required. \\
    Stan MCMC & Initial pilot consistent with Laplace point metrics & Full diagnostics and posterior predictive checks required. \\
    LANDFIRE time matching & LF2024 used for 2017--2023 panel & Time-matched LANDFIRE sensitivity required. \\
    \bottomrule
  \end{tabular}
\end{table}

\begin{figure}[!htbp]
  \centering
  \includegraphics[width=0.82\textwidth]{phase9/phase9_20260707_113214_metric_difference_bootstrap.png}
  \caption{Preliminary paired uncertainty check from segment-cluster bootstrap. This does not replace spatial block validation.}
  \label{fig:bootstrap}
\end{figure}
\FloatBarrier

\section{Discussion}

The main result is bounded. Public data can support a reproducible,
screening-level framework for retrospective external wildfire exposure ranking
of transmission-line receptor assets. The output is a ranked exposure screen,
not a statement about ignition causality, equipment condition, outage,
damage, or responsibility.

The temporal-holdout results do not show broad Bayesian classification
superiority. L2 logistic and Bayesian Laplace have nearly identical ROC-AUC,
PR-AUC, Brier score, and top-decile capture. This is consistent with the
relationship between Gaussian-prior logistic regression and L2 penalization.
The Bayesian contribution is instead decision-facing: posterior draws convert a
fitted exposure model into posterior exceedance probability and posterior
top-decile rank probability. These quantities are better aligned with
rare-event triage than a single deterministic score, because they express
uncertainty in threshold and rank decisions.

The deterministic comparator remains useful because it is transparent and easy
to audit. Its limitation is that it gives one ranking without posterior rank
uncertainty. Flexible tree-based ML baselines do not improve this temporal
holdout, but this result should be interpreted cautiously. It supports careful
feature design and temporal validation, not a general rejection of ML.

Existing wildfire risk tools target different questions. Powerline ignition and
equipment-failure models generally require private utility data and address
grid-to-fire causality. Fire-spread simulators can represent richer physical
dynamics but require scenario assumptions and are not automatically calibrated
for public-data segment ranking. The present framework occupies a narrower
space: public-data fire-to-grid external exposure ranking with a posterior
decision layer.

\section{Limitations}

Several limitations are important for a preprint. First, the pipeline estimates
external exposure only. It does not estimate ignition causality, equipment
failure, outage probability, damage, or legal responsibility. Second, the CAL
FIRE cause field is a perimeter-level attribute; electrical-power overlap is
therefore only a diagnostic label. Third, LF2024 static fuel layers may be
temporally mismatched with 2017--2023 exposure years; time-matched LANDFIRE
sensitivity is required. Fourth, FBFM40 must be encoded categorically or
removed from ablations. Fifth, the current label uses a fixed 1 km buffer;
250 m, 500 m, 1000 m, and 2000 m buffer sensitivities should be reported.
Sixth, temporal holdout is not spatial block validation. Finally, the Stan run
is a pilot; full diagnostics, longer chains, convergence summaries, and
posterior predictive checks are still needed.

\section{Conclusion}

This paper presents a public-data framework for retrospective external wildfire
exposure ranking of transmission-line receptor assets. The workflow builds a
2017--2023 segment-year panel, uses cause-screened exposure labels, and compares
deterministic, Bayesian, and machine-learning baselines under a 2022--2023
temporal holdout. Bayesian and L2 logistic point-prediction performance is
similar. The Bayesian formulation adds posterior exceedance and posterior
top-decile rank probabilities, with posterior rank probability emphasized as
the main decision quantity under rare-event calibration limits. The framework
should not be interpreted as a powerline ignition, equipment-failure, outage,
damage, or responsibility model.

\appendix

\section{Covariate Diagnostics}

The following maps document the public covariate layers used to construct
dynamic weather, terrain, and fuel-structure predictors. They are included as
quality-control and covariate-diagnostic figures rather than as separate
scientific claims.

\begin{figure}[!htbp]
  \centering
  \includegraphics[width=0.92\textwidth]{data_previews/phase2_gridmet_preview_2021.png}
  \caption{gridMET annual fire-weather covariate layers for 2021.}
  \label{fig:gridmet-layers}
\end{figure}

\begin{figure}[!htbp]
  \centering
  \includegraphics[width=0.86\textwidth]{data_previews/phase2_dem_preview.png}
  \caption{USGS 3DEP terrain covariate layers.}
  \label{fig:dem-layers}
\end{figure}

\begin{figure}[!htbp]
  \centering
  \includegraphics[width=0.92\textwidth]{data_previews/phase2_landfire_preview.png}
  \caption{LANDFIRE LF2024 fuel and canopy covariate layers.}
  \label{fig:landfire-layers}
\end{figure}

\begin{figure}[!htbp]
  \centering
  \includegraphics[width=0.92\textwidth]{phase3/phase3_covariate_qc_distributions.png}
  \caption{Covariate distribution diagnostics.}
  \label{fig:covariate-distributions}
\end{figure}

\begin{figure}[!htbp]
  \centering
  \includegraphics[width=0.92\textwidth]{phase3/phase3_spatial_qc_maps.png}
  \caption{Spatial covariate quality-control maps.}
  \label{fig:spatial-qc}
\end{figure}

\section{Reproducibility Notes}

All data processing, feature construction, model fitting, validation, and figure
generation are scripted in the project repository. Public-source citations are
still marked as TODO placeholders until a bibliography is added.

\end{document}
